\documentclass[10pt,a4paper]{amsart}
\usepackage{amsmath,amssymb,amsthm,graphicx}
\usepackage{verbatim}
\newenvironment{reference}{\comment}{\endcomment}
%\newenvironment{reference}{}{}
\newenvironment{slogan}{\comment}{\endcomment}
\newenvironment{history}{\comment}{\endcomment}

% For commutative diagrams we use Xy-pic
\usepackage[all]{xy}

% We use 2cell for 2-commutative diagrams.
\xyoption{2cell}
% Theorem environments.
%
\theoremstyle{plain}
\newtheorem{theorem}[subsection]{Theorem}
\newtheorem{proposition}[subsection]{Proposition}
\newtheorem{lemma}[subsection]{Lemma}

\theoremstyle{definition}
\newtheorem{definition}[subsection]{Definition}
\newtheorem{example}[subsection]{Example}
\newtheorem{exercise}[subsection]{Exercise}
\newtheorem{situation}[subsection]{Situation}

\theoremstyle{remark}
\newtheorem{remark}[subsection]{Remark}
\newtheorem{remarks}[subsection]{Remarks}

\numberwithin{equation}{subsection}

% Macros
%
\def\lim{\mathop{\mathrm{lim}}\nolimits}
\def\colim{\mathop{\mathrm{colim}}\nolimits}
\def\Spec{\mathop{\mathrm{Spec}}}
\def\Hom{\mathop{\mathrm{Hom}}\nolimits}
\def\Ext{\mathop{\mathrm{Ext}}\nolimits}
\def\SheafHom{\mathop{\mathcal{H}\!\mathit{om}}\nolimits}
\def\SheafExt{\mathop{\mathcal{E}\!\mathit{xt}}\nolimits}
\def\Sch{\mathit{Sch}}
\def\Mor{\mathop{\mathrm{Mor}}\nolimits}
\def\Ob{\mathop{\mathrm{Ob}}\nolimits}
\def\Sh{\mathop{\mathit{Sh}}\nolimits}
\def\NL{\mathop{N\!L}\nolimits}
\def\CH{\mathop{\mathrm{CH}}\nolimits}
\def\proetale{{pro\text{-}\acute{e}tale}}
\def\etale{{\acute{e}tale}}
\def\QCoh{\mathit{QCoh}}
\def\Ker{\mathop{\mathrm{Ker}}}
\def\Im{\mathop{\mathrm{Im}}}
\def\Coker{\mathop{\mathrm{Coker}}}
\def\Coim{\mathop{\mathrm{Coim}}}

% Boxtimes
%
\DeclareMathSymbol{\boxtimes}{\mathbin}{AMSa}{"02}

%
% Macros for moduli stacks/spaces
%
\def\QCohstack{\mathcal{QC}\!\mathit{oh}}
\def\Cohstack{\mathcal{C}\!\mathit{oh}}
\def\Spacesstack{\mathcal{S}\!\mathit{paces}}
\def\Quotfunctor{\mathrm{Quot}}
\def\Hilbfunctor{\mathrm{Hilb}}
\def\Curvesstack{\mathcal{C}\!\mathit{urves}}
\def\Polarizedstack{\mathcal{P}\!\mathit{olarized}}
\def\Complexesstack{\mathcal{C}\!\mathit{omplexes}}
% \Pic is the operator that assigns to X its picard group, usage \Pic(X)
% \Picardstack_{X/B} denotes the Picard stack of X over B
% \Picardfunctor_{X/B} denotes the Picard functor of X over B
\def\Pic{\mathop{\mathrm{Pic}}\nolimits}
\def\Picardstack{\mathcal{P}\!\mathit{ic}}
\def\Picardfunctor{\mathrm{Pic}}
\def\Deformationcategory{\mathcal{D}\!\mathit{ef}}
\usepackage{cleveref}
\crefname{originallemma}{original lemma}{original lemmas}
\Crefname{originallemma}{Original lemma}{Original lemmas}
\makeatletter
\newlabel{lemma-internal-hom-fp-into-qcoh}{{0GQN}{source}{}{sourcepointer-1}{}}
\newlabel{lemma-internal-hom-fp-into-qcoh@cref}{{[originallemma][0][]0GQN}{[1][0][]source}{}{}{}}
\makeatother
\crefname{originalreference}{original reference}{original references}
\Crefname{originalreference}{Original reference}{Original references}
\makeatletter
\newlabel{sites-modules-lemma-pullback-internal-hom}{{0GN2}{source}{}{sourcepointer-2}{}}
\newlabel{sites-modules-lemma-pullback-internal-hom@cref}{{[originalreference][0][]0GN2}{[1][0][]source}{}{}{}}
\makeatother
\crefname{originalreference}{original reference}{original references}
\Crefname{originalreference}{Original reference}{Original references}
\makeatletter
\newlabel{stacks-sheaves-remark-flat}{{06TX}{source}{}{sourcepointer-3}{}}
\newlabel{stacks-sheaves-remark-flat@cref}{{[originalreference][0][]06TX}{[1][0][]source}{}{}{}}
\makeatother
\crefname{originallemma}{original lemma}{original lemmas}
\Crefname{originallemma}{Original lemma}{Original lemmas}
\makeatletter
\newlabel{lemma-flat-morphism-and-Q}{{0GQJ}{source}{}{sourcepointer-4}{}}
\newlabel{lemma-flat-morphism-and-Q@cref}{{[originallemma][0][]0GQJ}{[1][0][]source}{}{}{}}
\makeatother
\title{Internal hom commutes with flat base change in QCoh of algebraic stacks}
\author{The Stacks Project Authors}
\date{stacks-project:stacks-cohomology.tex:local-snapshot-20261012; GNU FDL 1.2 or later}
\begin{document}
\maketitle
\noindent\textit{Editorial note.} All delivered mathematical text is the author's own, copied verbatim from the pinned Stacks Project chapter file stacks-cohomology.tex; no statement or proof step was rewritten, reordered or removed.\par\medskip
\noindent\textit{Editorial note.} The counted span cites the lemmas recorded in assembly.external\_references with canonical tags as navigation pointers.\par\medskip
\noindent\textit{Allowed standard prerequisites (editorial).}\begin{itemize}
\item Pullback and internal hom on ringed sites. Stacks Project, Modules on Sites, Lemma modules-lemma-flat-pushforward-pullback-hom (internal hom pullback identity). Gives the sheaf identity \$f\textasciicircum{}*\textbackslash{}mathcal\{H\}om(\textbackslash{}mathcal\{G\},\textbackslash{}mathcal\{H\}) = \textbackslash{}mathcal\{H\}om(f\textasciicircum{}*\textbackslash{}mathcal\{G\}, f\textasciicircum{}*\textbackslash{}mathcal\{H\})\$.
\item Flat morphisms and hom in QCoh\textasciicircum{}\{fbc\}. Stacks Project, Cohomology of Algebraic Stacks, Section on flat morphisms, Lemma stacks-cohomology-lemma-flat-morphism-and-Q. Identifies hom of the qc category with sheaf Hom under flat base change.
\end{itemize}\medskip
\ifdefined\hypertarget\hypertarget{sourcepointer-1}{}\fi
\noindent\textit{Source reference.} Original lemma 0GQN: sources/stacks-chow/tags.txt, source lines 20868--20868. Cross-chapter standard prerequisite cited by the author inside the counted span; it is used as a pointer and its own proof is not part of this item.\par\medskip
\ifdefined\hypertarget\hypertarget{sourcepointer-2}{}\fi
\noindent\textit{Source reference.} Original reference 0GN2: sources/stacks-chow/tags.txt, source lines 20777--20777. Cross-chapter standard prerequisite cited by the author inside the counted span; it is used as a pointer and its own proof is not part of this item.\par\medskip
\ifdefined\hypertarget\hypertarget{sourcepointer-3}{}\fi
\noindent\textit{Source reference.} Original reference 06TX: sources/stacks-chow/tags.txt, source lines 8726--8726. Cross-chapter standard prerequisite cited by the author inside the counted span; it is used as a pointer and its own proof is not part of this item.\par\medskip
\ifdefined\hypertarget\hypertarget{sourcepointer-4}{}\fi
\noindent\textit{Source reference.} Original lemma 0GQJ: sources/stacks-chow/tags.txt, source lines 20864--20864. Cross-chapter standard prerequisite cited by the author inside the counted span; it is used as a pointer and its own proof is not part of this item.\par\medskip
\begin{lemma}
\label{lemma-flat-morphism-and-hom}
Let $f : \mathcal{X} \to \mathcal{Y}$ be a flat morphism of algebraic stacks.
Let $\mathcal{F}$ be an
$\mathcal{O}_\mathcal{Y}$-module of finite presentation and let
$\mathcal{G}$ be a quasi-coherent $\mathcal{O}_\mathcal{Y}$-module.
Then $f^*hom(\mathcal{F}, \mathcal{G}) = hom(f^*\mathcal{F}, f^*\mathcal{G})$
with notation as in Lemma \ref{lemma-internal-hom-fp-into-qcoh}.
\end{lemma}

\begin{proof}
We have $f^*\SheafHom_{\mathcal{O}_\mathcal{Y}}(\mathcal{F}, \mathcal{G}) =
\SheafHom_{\mathcal{O}_\mathcal{X}}(f^*\mathcal{F}, f^*\mathcal{G})$
by Modules on Sites, Lemma \ref{sites-modules-lemma-pullback-internal-hom}.
(Observe that this step is not where the flatness of $f$ is used
as the morphism of ringed topoi associated to $f$ is always flat, see
Sheaves on Stacks, Remark \ref{stacks-sheaves-remark-flat}.)
Then apply Lemma \ref{lemma-flat-morphism-and-Q} (and here we do use
flatness of $f$).
\end{proof}

\end{document}
